% SGA 3，第 VII_B 篇，英文重建稿——第 1 轮。
% 范围：权威底本分篇第 77 页，第 4.3.1 节全文，止于第 4.3.2 节之前。
% 合并阅读本页码： 590.
% 再版印刷页码： 579.
% 本范围内原版页边定位： 546--547.
% 本范围内编者注： 137.
% 权威底本： Polo--Gille Exp7B-13oct24.pdf,
% SHA-256 D303690E0195654BBE7C3DAA5D89EA23E24CD709A3575217234AC19AE8836694.
% 曾参考 Jacob Reinhold 的 CC BY 4.0 Markdown 作为比对谱系。
% 本组件没有图表；全部数学均为原生 LaTeX。

% 2026-10-22：对照英文底本逐段复核；保留原有译文，核对公式、引用键、脚注与原生图路径。
\subsubsection*{4.3.1}
\phantomsection
\label{sga3:VIIB:4.3.1}

沿用 \SGARefLink{sga3:VIIB:4.3}{4.3} 的记号，来确定满足
\[
  \varepsilon_U(\xi)=1,
  \qquad
  \Delta_U(\xi)=\xi\otimes\xi.
\]
的元素 \(\xi\in U=U_p(L)\)。每个元素 \(\xi\in U\) 都有唯一的表示
\[
  \xi
  =
  \sum_{n\in\mathcal P}
  \xi_n\frac{x^n}{n!},
  \qquad
  \xi_n\in C.
\]
条件 \(\varepsilon_U(\xi)=1\) 给出
\[
  \xi_0=1,
\]
其中 \(0\) 表示 \(\mathcal P\) 中所有分量都为零的元素。条件
\(\Delta_{U_p(L)}(\xi)=\xi\otimes\xi\) 给出
\[
  \sum_{\substack{n\in\mathcal P\\m\leq n}}
  \xi_n
  \frac{x^m}{m!}
  \otimes
  \frac{x^{\,n-m}}{(n-m)!}
  =
  \sum_{q,r\in\mathcal P}
  \xi_q\xi_r
  \frac{x^q}{q!}
  \otimes
  \frac{x^r}{r!}.
\]
等价地，
\[
  \xi_{q+r}=\xi_q\xi_r
  \quad\text{若 }q_i+r_i<p\text{ 对每个 }i,
\]
以及
\[
  \xi_q\xi_r=0
  \quad\text{其他情形}.
\]

令 \(\xi_i\) 表示满足 \(n_i=1\) 且对 \(j\ne i\) 有 \(n_j=0\) 的系数 \(\xi_n\)。上述条件恰好表明
\[
  \xi_n=\prod_i\xi_i^{\,n_i}
  \quad(n\in\mathcal P),
  \qquad
  \xi_i^p=0
  \quad(i\in I).
\]
由此得到：

\medskip
\noindent\textbf{命题。}
\phantomsection
\label{sga3:VIIB:4.3.1:group-like-parametrization}
{\itshape 设 \(k\) 为一个特征为 \(p>0\) 的局部伪紧环\footnote{%
\label{sga3:VIIB:nde:137}编者增加了 \(k\) 为局部环的假设，以使每个拓扑平坦的伪紧 \(k\)-模都是拓扑自由的；见证明。}，\(\mathbf L\) 为 \(\mathbf O_k\) 上的一个平坦 \(p\)-Lie代数，\(C\) 为
\(\operatorname{Alf}/k\) 的一个对象。记
\[
  \sqrt[p]{0_C}
  =
  \{x\in C\mid x^p=0\}.
\]
则存在一个“关于 \(C\) 函子性”的双射
\[
  \mathbf L(C)\otimes_C\sqrt[p]{0_C}
  \xrightarrow{\ \sim\ }
  \mathscr G_p(\mathbf L)(C).
\]}

\textit{证明。}
由 \SGARefLink{sga3:VIIB:1.2.3.F}{注~1.2.3.F}，可以选取
\(\mathbf L(C)\) 的一组 \(C\)-基
\[
  ({}^C x_i)_{i\in I}
\]
使得对 \(\operatorname{Alf}/k\) 中的每个态射 \(\phi:C\to D\)，映射
\(\mathbf L(\phi)\) 都将 \({}^C x_i\) 送至 \({}^D x_i\)。前面的计算表明
\[
  \sum_i({}^C x_i)\otimes\xi_i
  \longmapsto
  \sum_{n\in\mathcal P}
  \left(\prod_i\xi_i^{\,n_i}\right)
  \frac{({}^C x)^n}{n!}
\]
给出从 \(\mathbf L(C)\otimes_C\sqrt[p]{0_C}\) 到
\(\mathscr G_p(\mathbf L)(C)\) 的一个函子性双射。
